\section{The fault menu}
\label{sec:faultmenu}

The stimulus side of the benchmark is a catalogue of parametric fault generators.
Each generator emits a true-error time series---one value per epoch, the hidden
$e(k)$ the scorer of Section~\ref{sec:method} compares against the monitor's
declared protection level---together with a \code{FaultSpec} recording the fault's
class, its worst-case injected offset, and, crucially, its \emph{detectability}.
The generators are \Modelled{}: they are representative, parametric, and
deliberately simple, not measured captures and not an exhaustive adversary. The
detectability partition, by contrast, is \Cited{}: which faults can be detected in
principle is a settled impossibility question, not a modelling choice.

\subsection{The ten fault classes}
\label{sec:classes}

Table~\ref{tab:faults} lists the catalogue and each class's offset profile and
detectability, as implemented in \code{src/benchmark/faults.rs}. The offset
profiles are elementary functions of the epoch index $k$; what matters is not their
sophistication but that they span the qualitatively distinct ways a timing service
loses integrity---slow domain drift, an abrupt slam, free-running holdover growth,
constant and ramping delays, the two provably undetectable delays, a
half-round-trip-time asymmetry, and the two partial-compromise geometries
(path-selective and $k$-of-$N$ quorum).

\begin{table}[t]
  \centering
  \caption{The fault menu (\Modelled{} generators; \Cited{} detectability). Offset
  profiles are functions of the epoch index $k$; $\mathrm{rate}$, $\mathrm{step}$,
  $d_{\mathrm{aging}}$, $r_{\mathrm{rw}}$, and the delay magnitudes are
  per-scenario parameters. The two undetectable classes are fixed by Mizrahi
  \citep{mizrahi7384}; the asymmetric-delay offset is fixed at half the path
  asymmetry by Narula and Humphreys \citep{narula2018}.}
  \label{tab:faults}
  \begin{tabular}{l l l}
    \toprule
    class & true-error profile $e(k)$ & detectability \\
    \midrule
    domain divergence     & $\mathrm{rate}\cdot k$ (slow drift)             & detectable \\
    clock slam            & $\mathrm{step}\cdot\mathbf 1[k\ge k_0]$          & detectable \\
    holdover coast        & $\tfrac12 d_{\mathrm{aging}}k^{2}+r_{\mathrm{rw}}k$ & detectable \\
    static delay          & constant offset                                 & detectable \\
    incremental delay     & $\mathrm{rate}\cdot k$ (ramp)                    & detectable \\
    symmetric delay       & constant offset                                 & \textbf{undetectable} \\
    asymmetric delay      & $\tfrac12\cdot(\text{path asymmetry})$           & detectable \\
    replay within freshness & constant offset                               & \textbf{undetectable} \\
    path-selective        & $\text{offset}\cdot(\text{fraction affected})$  & detectable \\
    $k$-of-$N$ quorum     & $\text{offset}\cdot(k/N)$                        & detectable \\
    \bottomrule
  \end{tabular}
\end{table}

The holdover-coast profile is monotone and quadratic-plus-linear in $k$
($\tfrac12 d_{\mathrm{aging}}k^{2}+r_{\mathrm{rw}}k$), matching the deterministic
shape of an aging, random-walk-driven free-running clock; the benchmark's own unit
test asserts this series is non-decreasing. The path-selective and quorum classes
model \emph{partial} compromise: a delay applied to a fraction of paths, or $k$ of
$N$ sources corrupted together, produces an effective fused offset scaled by that
fraction ($\text{fraction affected}$, or $k/N$)---detectable while an uncorrupted
quorum remains. These are representative geometries, not a claim about any specific
fusion rule.

\subsection{The detectability partition is \Cited{}, not modelled}
\label{sec:detectability}

Each generator carries a \code{Detectability} flag with two values,
\code{Detectable} and \code{Undetectable}, and the assignment is a citation, not a
design decision.

\paragraph{The undetectable set (Mizrahi).}
A \emph{symmetric single-path delay}---an adversary that delays the timing signal
by the same amount in both directions of a two-way exchange---is provably
undetectable by any protocol, cryptographic or otherwise, because it is
indistinguishable from a genuine change in path length \citep{mizrahi7384}.
Cryptographic authentication does not help: the delayed message is authentic, only
late, and the lateness is symmetric so the two-way estimate is unbiased about a
\emph{wrong} value. The same holds for a \emph{replay inside the freshness
window}: a message replayed before its freshness bound expires passes every
freshness and authentication check by construction. Both classes are therefore
flagged \code{Undetectable}, and the consequence---elaborated in
Section~\ref{sec:verdict}---is that a monitor's only honest defence is to
\emph{carry enough protection level to absorb} the offset, never to claim
detection. The benchmark's \code{full\_menu} contains \emph{exactly two}
undetectable scenarios, and this count is asserted by a unit test, so the
undetectable set cannot silently grow or vanish.

\paragraph{The boundable boundary (Narula and Humphreys).}
The complementary boundary is the \emph{asymmetric} delay. When the induced path
asymmetry is $\Delta$, the maximum offset a two-way exchange cannot resolve is
$\Delta/2$---half the round-trip-time asymmetry \citep{narula2018}. The generator
encodes exactly this: the injected offset of an asymmetric-delay scenario is
$\tfrac12\cdot(\text{path asymmetry})$. We are precise about what this partition
asserts: a pure path asymmetry is not \emph{directly detected} from a single
two-way exchange---nothing in one exchange isolates the asymmetry itself. What is
true is that its induced offset is \emph{boundable (resolvable)}: an asymmetry,
unlike a symmetric delay, perturbs the two-way residual and can in principle be
bounded. The class therefore carries the code's \code{Detectable} flag as a
\emph{routing} decision---the scenario is sent to the coverage verdict
(Section~\ref{sec:coverage}) rather than to the absorption verdict reserved for the
provably undetectable---not as a claim that the asymmetry is pinpointed. The unit
test checks both facts: that the asymmetric offset is half the asymmetry, and that
it lies on the boundable (\code{Detectable}) side of the partition.

\subsection{Scope of the menu}
\label{sec:menuscope}

The menu is \Modelled{} and we bound it plainly. The generators are deterministic
representative profiles, not stochastic captures of real network or clock
behaviour; the offset magnitudes are illustrative parameters, not operational
tolerances; and the ten classes, while spanning the qualitatively distinct failure
modes, are not an exhaustive adversary catalogue. Deep-tail integrity rates such as
\num{e-7} cannot be measured on the short scenarios a demonstration runs---resolving
them would need far more epochs than any illustrative menu carries, a point we
return to in Section~\ref{sec:discussion}. What the menu \emph{does} guarantee is
that the honesty-critical structure---the two undetectable classes and the
half-asymmetry boundary---is fixed by citation and locked by test, so that no
monitor can be scored against a menu that quietly omits the faults it cannot
survive.
